\documentclass[11pt]{article}

\usepackage{amsmath,amssymb,amsthm}
\usepackage{hyperref}
\usepackage{booktabs}
\usepackage[margin=1in]{geometry}

\newtheorem{theorem}{Theorem}[section]
\newtheorem{lemma}[theorem]{Lemma}
\newtheorem{definition}[theorem]{Definition}
\newtheorem{corollary}[theorem]{Corollary}

\title{Generalized Open Maps in Lean~4:\\One Abstract Module, Five Bisimulation Theories,\\and a Bridge to OSLF}
\author{Zar \and Oru\v{z}i (AI-assisted collaboration credit)}
\date{March 2026}

\begin{document}

\maketitle

\begin{abstract}
Joyal, Nielsen, and Winskel showed in 1996 that open maps provide a
categorical framework unifying diverse notions of bisimulation.  Despite
its theoretical elegance, no machine-checked formalization of this
framework---let alone one instantiated to multiple concrete
calculi---has appeared in the literature.  We present a 138-line Lean~4
module, \texttt{BisimulationKit}, that captures the generalized open-map
pattern and instantiates it to five concrete bisimulation theories:
(1)~weak bisimilarity for the $\pi$/$\rho$-calculus,
(2)~branching bisimilarity in the sense of van Glabbeek,
(3)~quantale-weighted bisimilarity,
(4)~OSLF observational equivalence with a full Hennessy--Milner
characterization under image-finiteness, and
(5)~a MeTTa-to-$\rho$ shared-fragment bridge with self-bisimulation
witness.
The entire development is sorry-free, uses zero custom axioms, and
builds with Lean~4.28 and Mathlib.  We further provide 8 regression
theorems pinning the core equivalences against silent breakage.
\end{abstract}


%% ====================================================================
\section{Introduction}\label{sec:intro}
%% ====================================================================

The categorical framework of \emph{open maps} introduced by Joyal,
Nielsen, and Winskel~\cite{JoyalNielsenWinskel1996} provides a uniform
account of bisimulation: a morphism is open when it has the path-lifting
property, and two objects are bisimilar when they are connected by a span
of open maps.  This single definition recovers, as special cases, strong
bisimulation, weak bisimulation, branching bisimulation, and hereditary
history-preserving bisimulation, simply by varying the category of path
objects.  The framework has been influential in concurrency theory, yet
to our knowledge no machine-checked formalization has appeared---not in
Coq, Isabelle/HOL, Agda, or Lean.

The closest mechanized work is Pous's Coq library for coinduction
up-to~\cite{Pous2016}, which addresses \emph{proof techniques} for
bisimulation (enhanced coinduction, up-to-context) rather than the
categorical open-map \emph{definition} itself.  Sangiorgi and
Walker~\cite{SangiorgiWalker2001} provide the standard textbook
treatment of bisimulation for the $\pi$-calculus, but their definitions
are not cast in the open-map framework.  Van Glabbeek's
spectrum~\cite{vanGlabbeek1993,vanGlabbeek2001} classifies behavioral
equivalences by discriminating power; the open-map viewpoint shows these
arise from a single parametric construction.

\paragraph{Contribution.}
We present a Lean~4 development consisting of:
\begin{enumerate}
\item A 138-line abstract module \texttt{GeneralizedOpenMaps.lean}
  defining the \texttt{BisimulationKit} structure and its derived
  notions (path witnesses, span witnesses, strong variants).
\item Four concrete instantiations to weak $\pi$/$\rho$-bisimilarity
  (78~lines), branching bisimilarity (61~lines), quantale-weighted
  bisimilarity (51~lines), and OSLF observational equivalence
  (108~lines).
\item A 273-line distinction-graph module establishing the
  Hennessy--Milner characterization theorem under image-finiteness,
  with a 363-line unification layer connecting behavioral, evidence, and
  totality semantics.
\item Eight regression theorems (80~lines total) pinning the exported
  equivalences against refactoring breakage.
\end{enumerate}
The entire development is sorry-free, uses only the three standard Lean
axioms (\texttt{propext}, \texttt{Classical.choice}, \texttt{Quot.sound}),
and builds with Lean~4.28~\cite{Moura2021} against Mathlib~\cite{mathlib2020}.


%% ====================================================================
\section{The BisimulationKit Core}\label{sec:core}
%% ====================================================================

The central abstraction lives in the file
\texttt{Mettapedia/CategoryTheory/GeneralizedOpenMaps.lean} (138~lines).
We describe its components in the order they appear.

\begin{definition}[BisimulationKit]\label{def:kit}
A \texttt{BisimulationKit} over a state type $\alpha$ and an observable
type \texttt{Obs} is a record:
\begin{verbatim}
  structure BisimulationKit (α : Type u) (Obs : Type v) where
    step      : α → α → Prop
    stepStar  : α → α → Prop
    step_sub_star : ∀ {x y}, step x y → stepStar x y
    observable : α → Obs → Prop
\end{verbatim}
The field \texttt{step} is the one-step transition relation,
\texttt{stepStar} is its saturated (weak/reflexive-transitive) closure,
\texttt{step\_sub\_star} witnesses the embedding, and
\texttt{observable} is a family of observable predicates indexed by
\texttt{Obs}.
\end{definition}

Given a kit $K$, the module defines several derived notions.

\begin{definition}[Run and Extension]
A \texttt{Run K} is a source/target pair $(s, d)$ together with a proof
that $K.\mathit{stepStar}\; s\; d$.  An \texttt{Extension} of run $r$
to run $r'$ holds when they share the same source and $r'$ reaches
further: $K.\mathit{stepStar}\; r.\mathit{dst}\; r'.\mathit{dst}$.
\end{definition}

\begin{definition}[GOpen, GOpenStrong, ObsPreserving]\label{def:gopen}
For a binary relation $R$ on $\alpha$:
\begin{itemize}
\item \texttt{GOpen K R} (weak open lifting): for all $p, q$ with
  $R\;p\;q$ and $p' $ with $K.\mathit{step}\;p\;p'$, there exists $q'$
  with $K.\mathit{stepStar}\;q\;q'$ and $R\;p'\;q'$.
\item \texttt{GOpenStrong K R} (strong open lifting): the same, but
  requiring $K.\mathit{step}\;q\;q'$ instead of
  $K.\mathit{stepStar}\;q\;q'$.
\item \texttt{ObsPreserving K R}: for all $p, q$ with $R\;p\;q$ and
  observable $o$, $K.\mathit{observable}\;p\;o$ implies
  $K.\mathit{observable}\;q\;o$.
\end{itemize}
\end{definition}

The weak/strong distinction is precisely the difference between weak and
strong bisimulation in the standard process-algebraic sense.

\begin{definition}[Path Witnesses]
A \texttt{PathWitness K} (lines 55--59) bundles a symmetric relation
$R$ with \texttt{GOpen K R} and \texttt{ObsPreserving K R}.
A \texttt{StrongPathWitness K} (lines 62--66) replaces \texttt{GOpen}
with \texttt{GOpenStrong}.
\end{definition}

\begin{definition}[Bisimilarity]
Two states are \texttt{PathBisim K p q} (line 69) when there exists a
\texttt{PathWitness} whose relation contains $(p, q)$.
Similarly for \texttt{StrongPathBisim K} (line 73).
\end{definition}

\begin{definition}[Span Witness]
A \texttt{GOpenSpanWitness K} (lines 78--82) is structurally identical
to \texttt{PathWitness} in this minimal core---it bundles a symmetric
relation with \texttt{GOpen} and \texttt{ObsPreserving}.  The type name
documents the span-of-open-maps reading.  The corresponding existence
wrapper is \texttt{ESBisimilar K p q} (line 85).
\end{definition}

The core module proves two key structural theorems:

\begin{theorem}[\texttt{pathBisim\_iff\_esBisimilar}, line 120]\label{thm:path-es}
For any kit $K$ and states $p, q$:
\[
  \texttt{PathBisim}\;K\;p\;q \iff \texttt{ESBisimilar}\;K\;p\;q.
\]
\end{theorem}
\begin{proof}
Each direction converts between \texttt{PathWitness} and
\texttt{GOpenSpanWitness} via the bijections \texttt{toSpan} and
\texttt{toPath}, which simply repackage the same fields.
\end{proof}

\begin{theorem}[\texttt{strongPathBisim\_implies\_pathBisim}, line 128]\label{thm:strong-weak}
For any kit $K$:
\[
  \texttt{StrongPathBisim}\;K\;p\;q \implies \texttt{PathBisim}\;K\;p\;q.
\]
\end{theorem}
\begin{proof}
Given a \texttt{StrongPathWitness} $w$, construct a
\texttt{PathWitness} with the same relation, symmetry, and observable
preservation.  The lifting property weakens: the strong step witness
$K.\mathit{step}\;q\;q'$ is embedded into
$K.\mathit{stepStar}\;q\;q'$ via \texttt{step\_sub\_star}.
\end{proof}

The module also proves auxiliary lemmas \texttt{PathWitness.open\_symm}
and \texttt{PathWitness.obs\_symm} showing that the flip of a path
witness relation is again open and observable-preserving.


%% ====================================================================
\section{Instantiation 1: Weak Bisimilarity for $\pi$/$\rho$}\label{sec:pirho}
%% ====================================================================

The file \texttt{PiCalculus/WeakBisimOpenMapBridge.lean} (78~lines)
bridges the existing rho-calculus formalization to the generalized
open-map framework.

\begin{definition}[ObsName]
For a finite set of names $N : \texttt{Finset String}$,
$\texttt{ObsName}\;N = \{x : \texttt{String} \mid x \in N\}$ is the
subtype of observable channel names.
\end{definition}

\begin{definition}[PiRhoInst]
The kit \texttt{PiRhoInst N} instantiates \texttt{BisimulationKit Pattern (ObsName N)} with:
\begin{itemize}
\item \texttt{step} = \texttt{Nonempty (Reduces p q)}, wrapping the
  rho-calculus reduction relation in \texttt{Nonempty} for Prop membership;
\item \texttt{stepStar} = \texttt{Nonempty (ReducesStar p q)}, the
  reflexive-transitive closure;
\item \texttt{observable} = \texttt{RhoObservableSC p (.fvar o.1)},
  structural-congruence-respecting barb observability.
\end{itemize}
A second kit \texttt{PiRhoDerivedInst N} does the same for the derived
reduction relation ($\Rightarrow_d$ and $\Rightarrow_d^*$).
\end{definition}

The main results are genuine characterization theorems---both directions
are proven:

\begin{theorem}[\texttt{weakRestrictedBisim\_iff\_pathBisim}, lines 44--59]\label{thm:weak-pirho}
\[
  \texttt{WeakRestrictedBisim}\;N\;p\;q \iff
  \texttt{PathBisim}\;(\texttt{PiRhoInst}\;N)\;p\;q.
\]
\end{theorem}

\begin{theorem}[\texttt{weakRestrictedBisimD\_iff\_pathBisim}, lines 61--76]\label{thm:weakD-pirho}
\[
  \texttt{WeakRestrictedBisimD}\;N\;p\;q \iff
  \texttt{PathBisim}\;(\texttt{PiRhoDerivedInst}\;N)\;p\;q.
\]
\end{theorem}

Both proofs proceed by unpacking the existing bisimulation definition
(which stores a symmetric relation, a forward simulation, and observable
preservation) and repackaging it as a \texttt{PathWitness}, and
conversely.  The forward direction extracts the existential witness from
the \texttt{WeakRestrictedBisim} record; the backward direction
reconstructs it from the \texttt{PathWitness} fields, using the subtype
projection $o.\mathit{val} : \texttt{String}$ and membership
$o.\mathit{property} : o.\mathit{val} \in N$ to bridge between the
universally-quantified and subtype-indexed observable formats.


%% ====================================================================
\section{Instantiation 2: Branching Bisimilarity}\label{sec:branching}
%% ====================================================================

The file \texttt{PiCalculus/BranchingBisim.lean} (61~lines) adds a thin
naming layer that connects the strong-path variant of the open-map
framework to van Glabbeek's branching bisimulation~\cite{vanGlabbeek1993,vanGlabbeek2001}.

\begin{definition}[Open-Map Naming]
Three abbreviations make the correspondence explicit:
\begin{itemize}
\item \texttt{WeakBisimOM K p q} $\equiv$ \texttt{PathBisim K p q};
\item \texttt{BranchingBisimOM K p q} $\equiv$ \texttt{StrongPathBisim K p q};
\item \texttt{StutteringBranchingBisimOM K p q} $\equiv$ \texttt{StrongPathBisim K p q}.
\end{itemize}
The identification of branching and stuttering-branching bisimulation is
an artifact of the open-map setting, where both reduce to the same
strong-path notion---the stuttering condition is automatically
satisfied by the path structure.
\end{definition}

\begin{theorem}[\texttt{branching\_implies\_weak}, line 41]\label{thm:branching-weak}
For any kit $K$ and states $p, q$:
\[
  \texttt{BranchingBisimOM}\;K\;p\;q \implies
  \texttt{WeakBisimOM}\;K\;p\;q.
\]
\end{theorem}
\begin{proof}
Direct application of
Theorem~\ref{thm:strong-weak}
(\texttt{strongPathBisim\_implies\_pathBisim}).
\end{proof}

The file also provides two concrete bridge theorems connecting back to
the existing $\pi$/$\rho$-calculus definitions:

\begin{theorem}[\texttt{weakRestrictedBisim\_of\_branchingOM}]\label{thm:branch-concrete1}
\[
  \texttt{BranchingBisimOM}\;(\texttt{PiRhoInst}\;N)\;p\;q \implies
  \texttt{WeakRestrictedBisim}\;N\;p\;q.
\]
\end{theorem}

\begin{theorem}[\texttt{weakRestrictedBisimD\_of\_branchingOM}]\label{thm:branch-concrete2}
\[
  \texttt{BranchingBisimOM}\;(\texttt{PiRhoDerivedInst}\;N)\;p\;q \implies
  \texttt{WeakRestrictedBisimD}\;N\;p\;q.
\]
\end{theorem}

Both are one-liners composing the branching-to-weak implication with the
characterization theorems from Section~\ref{sec:pirho}.


%% ====================================================================
\section{Instantiation 3: Quantale-Weighted Bisimilarity}\label{sec:weighted}
%% ====================================================================

The file \texttt{Logic/WeightedOpenMaps.lean} (51~lines) instantiates
the open-map kit for quantale-labeled transition systems (QLTS), as
defined in the companion module
\texttt{Logic/ModalQuantaleSemantics.lean} (520~lines).

\begin{definition}[QLTS Step Relations]
For a QLTS with state type $S$, action type \texttt{Act}, and weight
lattice $Q$ (a complete lattice):
\begin{itemize}
\item \texttt{qltsStep qlts s t}: there exists an action $a$ such that
  $\texttt{qlts.trans}\;s\;a\;t \neq \bot$;
\item \texttt{qltsStepStar qlts}: the reflexive-transitive closure
  \texttt{ReflTransGen} of \texttt{qltsStep}.
\end{itemize}
\end{definition}

\begin{definition}[\texttt{QLTSInst}]
The kit \texttt{QLTSInst qlts} instantiates
\texttt{BisimulationKit S Act} with \texttt{step} = \texttt{qltsStep},
\texttt{stepStar} = \texttt{qltsStepStar}, and observable $s\;a$ iff
there exists a successor with non-$\bot$ weight on action $a$.
\end{definition}

\begin{theorem}[\texttt{weightedBisim\_iff\_gopen\_span}, line 46]\label{thm:weighted}
For any QLTS and states $s, t$:
\[
  \texttt{WeightedBisim}\;\mathit{qlts}\;s\;t \iff
  \texttt{WeightedGOpenSpanBisim}\;\mathit{qlts}\;s\;t.
\]
\end{theorem}
\begin{proof}
Direct application of Theorem~\ref{thm:path-es}
(\texttt{pathBisim\_iff\_esBisimilar}) at the \texttt{QLTSInst}
instantiation.
\end{proof}

This instantiation connects the open-map framework to the 520-line
modal quantale semantics development, which provides a complete
lattice-valued interpretation of modal formulas over QLTS.  The
connection to Kozen's $\mu$-calculus~\cite{Kozen1983} is made precise
in \texttt{ModalQuantaleSemantics.lean}, where fixpoint operators are
interpreted via Knaster--Tarski.


%% ====================================================================
\section{Instantiation 4: OSLF Observational Equivalence}\label{sec:oslf}
%% ====================================================================

The file \texttt{Logic/OSLFOpenMapBridge.lean} (108~lines) bridges the
open-map framework to Operational Semantics in Logical Form
(OSLF)~\cite{MeredithStayOSLF2020,leanOSLF2025}.

\begin{definition}[\texttt{OSLFInst}]
Given a step relation $R : \texttt{Pat} \to \texttt{Pat} \to \texttt{Prop}$
and an atom semantics $I : \texttt{AtomSem}$, the kit
\texttt{OSLFInst R I} sets both \texttt{step} and \texttt{stepStar} to
$R$ (the OSLF framework treats each reduction step as already
saturated), with \texttt{observable p a} $\equiv$ $I\;a\;p$.
\end{definition}

\begin{definition}[\texttt{FullOpenWitness}]
A \texttt{FullOpenWitness R I} (lines 52--57) bundles a symmetric
relation with:
\begin{itemize}
\item \texttt{open\_fwd}: \texttt{GOpen (OSLFInst R I) rel}---forward
  open lifting for $R$;
\item \texttt{open\_rev}: \texttt{GOpen (OSLFInst ($\lambda$\,a\,b.\;R\,b\,a) I) rel}---open
  lifting for the reverse relation $R^{-1}$;
\item \texttt{atom}: \texttt{ObsPreserving (OSLFInst R I) rel}---atom
  semantics preservation.
\end{itemize}
The bidirectionality is essential: the OSLF formula language includes
both diamond ($\Diamond$, requiring forward simulation) and box
($\Box$, requiring backward simulation via $R^{-1}$).
\end{definition}

The module proves four theorems forming a cascade of increasingly strong
consequences:

\begin{theorem}[\texttt{pathBisim\_implies\_bisimilar}, lines 44--48]\label{thm:oslf-bisim}
\[
  \texttt{PathBisim}\;(\texttt{OSLFInst}\;R\;I)\;p\;q \implies
  \texttt{Bisimilar}\;R\;p\;q.
\]
\end{theorem}
\begin{proof}
Extract the \texttt{PathWitness} and construct a
\texttt{StepBisimulation} via a private lemma
\texttt{pathWitness\_stepBisimulation}, which derives the backward
simulation direction from the symmetry of the witness relation.
\end{proof}

\begin{theorem}[\texttt{fullOpenWitness\_implies\_obsEq}, lines 88--92]\label{thm:oslf-obseq}
A \texttt{FullOpenWitness} implies OSLF observational equivalence
(\texttt{OSLFObsEq}): related states satisfy the same OSLF formulas.
\end{theorem}

\begin{theorem}[\texttt{fullOpenWitness\_implies\_indistObs}, lines 95--99]\label{thm:oslf-indist}
A \texttt{FullOpenWitness} implies indistinguishability in the
distinction graph (\texttt{indistObs}).
\end{theorem}

\begin{theorem}[\texttt{fullOpenWitness\_not\_distinguished}, lines 101--106]\label{thm:oslf-nodist}
A \texttt{FullOpenWitness} precludes the existence of any separating
formula: $\neg\;\texttt{distinguished}\;R\;I\;p\;q$.
\end{theorem}


%% ====================================================================
\section{The Distinction Graph and Hennessy--Milner Converse}\label{sec:hm}
%% ====================================================================

The distinction graph and its Hennessy--Milner characterization are
developed in two files:
\texttt{Logic/OSLFDistinctionGraph.lean} (273~lines) and
\texttt{Logic/OSLFKSUnificationSketch.lean} (363~lines).

\subsection{Distinction Graph Structure}

\begin{definition}[\texttt{indistObs}, \texttt{distinguished}]
For a step relation $R$ and atom semantics $I$:
\begin{itemize}
\item $\texttt{indistObs}\;R\;I\;p\;q$ holds iff $p$ and $q$ satisfy
  the same OSLF formulas: $\forall\,\varphi,\;
  \texttt{sem}\;R\;I\;\varphi\;p \iff \texttt{sem}\;R\;I\;\varphi\;q$.
\item $\texttt{distinguished}\;R\;I\;p\;q \equiv
  \neg\;\texttt{indistObs}\;R\;I\;p\;q$: some formula separates them.
\end{itemize}
\end{definition}

The module proves that \texttt{indistObs} is an equivalence relation
(reflexive, symmetric, transitive), yielding a \texttt{Setoid} and a
quotient type \texttt{ObsClass R I}.  The distinction graph is
irreflexive and symmetric.

\subsection{Bisimulation Invariance}

\begin{theorem}[\texttt{bisimulation\_invariant\_sem}]
If $E$ is a \texttt{StepBisimulation} for both $R$ and $R^{-1}$, and
$E$ preserves atom semantics, then $E$-related states satisfy the same
OSLF formulas.  The proof proceeds by structural induction on the
formula $\varphi$.
\end{theorem}

\begin{corollary}[\texttt{fullBisim\_implies\_indist}]
\texttt{FullBisimilar R I p q} implies \texttt{indistObs R I p q}.
\end{corollary}

\subsection{Hennessy--Milner Converse}

The converse---indistinguishability implies bisimilarity---requires
image-finiteness, following Hennessy and Milner~\cite{HennessyMilner1985}.

\begin{theorem}[\texttt{indistObs\_is\_stepBisimulation}]\label{thm:hm-fwd}
Under forward image-finiteness ($\forall\,p,\;\{q \mid R\;p\;q\}$ is
finite), \texttt{indistObs R I} is a \texttt{StepBisimulation} for $R$.
\end{theorem}

\begin{theorem}[\texttt{indistObs\_is\_revStepBisimulation}]\label{thm:hm-rev}
Under backward image-finiteness ($\forall\,p,\;\{q \mid R\;q\;p\}$ is
finite), \texttt{indistObs R I} is a \texttt{StepBisimulation} for
$R^{-1}$.
\end{theorem}

The proof of Theorem~\ref{thm:hm-rev} is noteworthy: it uses the
formula $\neg\Box\neg\Psi$ (encoded as
$\texttt{.imp}\;(\texttt{.box}\;(\texttt{.imp}\;\Psi\;\texttt{.bot}))\;\texttt{.bot}$)
to express the reverse-direction diamond ``there exists an $R$-predecessor
satisfying $\Psi$.''  The formula $\Psi$ is constructed as a finite
conjunction over the finitely many predecessors, with each conjunct
being a separating formula obtained from the
\texttt{separator\_of\_not\_obsEq} lemma.

\begin{theorem}[\texttt{indist\_iff\_fullBisim\_imageFinite}]\label{thm:hm-iff}
Under both forward and backward image-finiteness:
\[
  \texttt{indistObs}\;R\;I\;p\;q \iff \texttt{FullBisimilar}\;R\;I\;p\;q.
\]
\end{theorem}

This is the central Hennessy--Milner characterization theorem:
observational equivalence and bisimilarity coincide for image-finite
systems.

\subsection{Quotient and Measurement}

\begin{theorem}[\texttt{measurement\_factors\_through\_obsClass}]\label{thm:factor}
Any valuation $\mu : \texttt{Pat} \to \mathbb{R}$ that respects
indistinguishability factors through the quotient:
there exists $\mu_Q : \texttt{ObsClass}\;R\;I \to \mathbb{R}$ with
$\mu_Q\;[p] = \mu\;p$.
\end{theorem}

\subsection{Three-Layer Unification}

The 363-line module \texttt{OSLFKSUnificationSketch.lean} connects three
layers of semantics:
\begin{enumerate}
\item \textbf{Meredith}~\cite{MeredithRadestock2005}: behavioral
  semantics via bisimulation invariance and the Hennessy--Milner
  converse;
\item \textbf{Stay/Baez}: evidence semantics with threshold atoms,
  checker soundness, evidence revision, and rewrite-rule preservation;
\item \textbf{Knuth/Skilling}: the totality gate---faithful
  scalarization exists only for total orders, while the evidence
  lattice is non-total (the imprecision gate).
\end{enumerate}


%% ====================================================================
\section{Instantiation 5: MeTTa Shared-Fragment Bridge}\label{sec:metta-bridge}
%% ====================================================================

The MeTTa-calculus~\cite{MeredithMeTTaCalculus2024} shares the
quote/drop/parallel subset of $\rho$-calculus.
\texttt{MeTTaCalculus/\allowbreak{}Interoperability.lean}
(113~lines) provides a partial translation
\texttt{toRhoSharedProc?}\ that maps this shared fragment into
$\rho$-calculus terms.  The translation succeeds on
\texttt{pZero}, quote/drop, and parallel bags, and fails
gracefully outside the translatable fragment.

For any successfully translated process, the file proves that the
$\rho$-image carries a canonical open-map path-bisimulation witness:

\begin{theorem}[\texttt{shared\_translation\_pathBisim\_self}]
If $\texttt{toRhoSharedProc?}\;p = \texttt{some}\;rp$, then
$\texttt{PathBisim}\;\texttt{RhoSharedOpenInst}\;rp\;rp$.
\end{theorem}

\noindent
A span-style corollary \texttt{shared\_translation\_esBisim\_self}
extracts the corresponding $(E,S)$-bisimilarity witness, and
\texttt{paper\_shared\_to\_rho\_openmap\_self} (\texttt{PaperMap.lean},
65~lines) provides the paper-mapped wrapper.

\paragraph{Scope and future work.}
This result establishes that the shared-fragment translation is
well-typed with respect to the open-map framework---the translated image
is self-bisimilar under the $\rho$-calculus \texttt{BisimulationKit}.
It is \emph{not yet} a full cross-language bisimulation theorem.  The
next step is to prove forward/backward simulation for translated
shared-fragment steps, then upgrade to a full
MeTTa$\leftrightarrow$$\rho$ bisimulation theorem.


%% ====================================================================
\section{Regression and Theorem Pinning}\label{sec:regression}
%% ====================================================================

Two regression files guard the exported API against silent breakage
during refactoring:

\begin{itemize}
\item \texttt{PiCalculus/OpenMapBridgeRegression.lean} (34~lines):
  4~theorems pinning the $\pi$/$\rho$ bridge equivalences.
\item \texttt{Logic/OpenMapBridgeRegression.lean} (46~lines):
  4~theorems pinning the weighted and OSLF bridge exports.
\end{itemize}

Each regression theorem restates a core result and proves it by
\texttt{simpa using} (for iff) or direct application (for implications)
of the original.  For example:

\begin{verbatim}
  theorem weakRestrictedBisim_pathBisim_regression
      (N : Finset String) (p q : Pattern) :
      WeakRestrictedBisim N p q <->
        PathBisim (PiRhoInst N) p q := by
    simpa using weakRestrictedBisim_iff_pathBisim N p q
\end{verbatim}

The methodology is simple but effective: if a refactoring changes the
statement or proof of a core theorem in a way that breaks
compatibility, the regression file will fail to compile before the
change can be merged.  This is especially valuable in a Lean~4 +
Mathlib setting where upstream API changes can silently alter the
meaning of definitions.


%% ====================================================================
\section{Related Work and Conclusion}\label{sec:conclusion}
%% ====================================================================

\paragraph{Joyal--Nielsen--Winskel (1996).}
The original categorical framework~\cite{JoyalNielsenWinskel1996}
establishes that open maps in presheaf categories yield a uniform
definition of bisimulation.  Our formalization captures the operational
core of this insight---the path-lifting property---without formalizing
the full presheaf-category infrastructure.

\paragraph{Van Glabbeek's spectrum.}
The linear-time--branching-time
spectrum~\cite{vanGlabbeek1993,vanGlabbeek2001} classifies behavioral
equivalences by discriminating power.  Our branching bisimulation
instantiation (Section~\ref{sec:branching}) connects to this spectrum
via the strong-path variant.

\paragraph{Pous (2016).}
Coinduction up-to in Coq~\cite{Pous2016} addresses proof techniques
for bisimulation (enhanced coinduction, up-to-context, up-to-equivalence)
rather than the categorical open-map definition.  The two approaches are
complementary: Pous's techniques could be layered on top of our
framework to streamline proofs of specific bisimulation results.

\paragraph{Milner and Sangiorgi--Walker.}
The standard $\pi$-calculus theory~\cite{Milner1999,SangiorgiWalker2001}
provides the definitions our instantiation targets.  The rho-calculus
extension~\cite{MeredithRadestock2005} adds reflection and name-quoting.

\paragraph{OSLF.}
Meredith and Stay's Operational Semantics in Logical
Form~\cite{MeredithStayOSLF2020} provides the formula language and
behavioral-equivalence definitions targeted by our OSLF bridge.
The Lean formalization~\cite{leanOSLF2025} provides the infrastructure
(\texttt{OSLFFormula}, \texttt{sem}, \texttt{StepBisimulation}) on which
Sections~\ref{sec:oslf}--\ref{sec:hm} build.

\paragraph{Summary.}
Table~\ref{tab:files} summarizes the formalization.  The key takeaway is
that a single 138-line abstract module, when instantiated with
appropriate step relations and observables, recovers four distinct
bisimulation theories with fully machine-checked proofs.  The
Hennessy--Milner iff (Theorem~\ref{thm:hm-iff}) closes the
logical characterization under image-finiteness, and the measurement
factoring theorem (Theorem~\ref{thm:factor}) connects the
behavioral theory to quantitative reasoning.

\begin{table}[t]
\centering
\caption{Formalization summary.}\label{tab:files}
\begin{tabular}{@{}llr@{}}
\toprule
\textbf{File} & \textbf{Role} & \textbf{Lines} \\
\midrule
\texttt{GeneralizedOpenMaps.lean} & Abstract core (BisimulationKit) & 138 \\
\texttt{WeakBisimOpenMapBridge.lean} & Weak $\pi$/$\rho$ bisim $\leftrightarrow$ PathBisim & 78 \\
\texttt{BranchingBisim.lean} & Branching bisim via StrongPathBisim & 61 \\
\texttt{WeightedOpenMaps.lean} & Quantale-weighted bisim & 51 \\
\texttt{OSLFOpenMapBridge.lean} & OSLF observational equivalence & 108 \\
\texttt{OSLFDistinctionGraph.lean} & Distinction graph + Hennessy--Milner & 273 \\
\texttt{OSLFKSUnificationSketch.lean} & Three-layer unification & 363 \\
\texttt{ModalQuantaleSemantics.lean} & Quantale modal logic & 520 \\
\texttt{Interoperability.lean} & MeTTa shared-fragment bridge & 113 \\
\texttt{OpenMapBridgeRegression.lean} (2 files) & Regression guards & 80 \\
\midrule
\textbf{Total} & & \textbf{1785} \\
\bottomrule
\end{tabular}
\end{table}

\paragraph{Future work.}
Natural next steps include: (i)~formalizing the full presheaf-category
machinery to connect our operational core to the original JNW categorical
framework; (ii)~adding up-to techniques \`{a} la Pous to streamline
bisimulation proofs; and (iii)~instantiating the kit for additional
calculi (CCS, CSP, ambient calculus) to further validate the
``one module, many theories'' claim.

\bibliographystyle{plain}
\bibliography{references}

\end{document}
